% Options for packages loaded elsewhere
\PassOptionsToPackage{unicode}{hyperref}
\PassOptionsToPackage{hyphens}{url}
\PassOptionsToPackage{dvipsnames,svgnames,x11names}{xcolor}
%
\documentclass[
  10pt,
  a4paper,
]{article}
\usepackage{amsmath,amssymb}
\usepackage{setspace}
\usepackage{iftex}
\ifPDFTeX
  \usepackage[T1]{fontenc}
  \usepackage[utf8]{inputenc}
  \usepackage{textcomp} % provide euro and other symbols
\else % if luatex or xetex
  \usepackage{unicode-math} % this also loads fontspec
  \defaultfontfeatures{Scale=MatchLowercase}
  \defaultfontfeatures[\rmfamily]{Ligatures=TeX,Scale=1}
\fi
\usepackage{lmodern}
\ifPDFTeX\else
  % xetex/luatex font selection
  \setmainfont[]{Liberation Serif}
  \setsansfont[]{Liberation Sans}
  \setmonofont[]{DejaVu Sans Mono}
\fi
% Use upquote if available, for straight quotes in verbatim environments
\IfFileExists{upquote.sty}{\usepackage{upquote}}{}
\IfFileExists{microtype.sty}{% use microtype if available
  \usepackage[]{microtype}
  \UseMicrotypeSet[protrusion]{basicmath} % disable protrusion for tt fonts
}{}
\makeatletter
\@ifundefined{KOMAClassName}{% if non-KOMA class
  \IfFileExists{parskip.sty}{%
    \usepackage{parskip}
  }{% else
    \setlength{\parindent}{0pt}
    \setlength{\parskip}{6pt plus 2pt minus 1pt}}
}{% if KOMA class
  \KOMAoptions{parskip=half}}
\makeatother
\usepackage{xcolor}
\usepackage[margin=22mm]{geometry}
\setlength{\emergencystretch}{3em} % prevent overfull lines
\providecommand{\tightlist}{%
  \setlength{\itemsep}{0pt}\setlength{\parskip}{0pt}}
\setcounter{secnumdepth}{-\maxdimen} % remove section numbering
\ifLuaTeX
\usepackage[bidi=basic]{babel}
\else
\usepackage[bidi=default]{babel}
\fi
\babelprovide[main,import]{korean}
\ifPDFTeX
\else
\babelfont{rm}[]{Liberation Serif}
\fi
% get rid of language-specific shorthands (see #6817):
\let\LanguageShortHands\languageshorthands
\def\languageshorthands#1{}
\usepackage{xeCJK}
\setCJKmainfont{Noto Serif CJK KR}
\setCJKsansfont{Noto Sans CJK KR}
\setCJKmonofont{Noto Sans Mono CJK KR}
\xeCJKsetup{CJKspace=true}
\usepackage{xurl}
\usepackage{fvextra}
\usepackage{fancyhdr}
\AtBeginDocument{\DefineVerbatimEnvironment{verbatim}{Verbatim}{breaklines=true,breakanywhere=true,fontsize=\small}}
\setlength{\emergencystretch}{4em}
\setlength{\headheight}{15pt}
\pagestyle{fancy}\fancyhf{}
\fancyhead[L]{\small 중성자별 구조와 비방사 진동}
\fancyfoot[C]{\thepage}
\renewcommand{\headrulewidth}{0.3pt}
\setcounter{tocdepth}{2}
\ifLuaTeX
  \usepackage{selnolig}  % disable illegal ligatures
\fi
\usepackage{bookmark}
\IfFileExists{xurl.sty}{\usepackage{xurl}}{} % add URL line breaks if available
\urlstyle{same}
\hypersetup{
  pdftitle={A21 · 같은 경계값 문제에서 작은 모드 차이를 분해하기},
  pdflang={ko-KR},
  colorlinks=true,
  linkcolor={Maroon},
  filecolor={Maroon},
  citecolor={Blue},
  urlcolor={Blue},
  pdfcreator={LaTeX via pandoc}}

\title{A21 · 같은 경계값 문제에서 작은 모드 차이를 분해하기}
\usepackage{etoolbox}
\makeatletter
\providecommand{\subtitle}[1]{% add subtitle to \maketitle
  \apptocmd{\@title}{\par {\large #1 \par}}{}{}
}
\makeatother
\subtitle{과제 큐시트}
\author{}
\date{}

\begin{document}
\maketitle

\setstretch{1.12}
\textbf{선행:} A19, A20　 \textbf{독립 작업:} 4--6시간 (계획치)

\subsection{같은 문제의 해라는 주장을 무엇으로
방어하는가?}\label{uxac19uxc740-uxbb38uxc81cuxc758-uxd574uxb77cuxb294-uxc8fcuxc7a5uxc744-uxbb34uxc5c7uxc73cuxb85c-uxbc29uxc5b4uxd558uxb294uxac00}

두 방향에서 얻은 벡터의 determinant가 작다. 이것만으로 다른 EOS·다른
질량에서 얻은 해도 연결할 수 있을까? 잔차 앞에 필요한 조건을 먼저
말하자.

먼저 한두 문장으로 답한 뒤 아래 설명을 읽자. 첫 예측이 틀렸거나 아직
모르겠다고 적는 것은 감점하지 않는다.

\subsection{필요한 관계와 읽을
위치}\label{uxd544uxc694uxd55c-uxad00uxacc4uxc640-uxc77duxc744-uxc704uxce58}

\(D=\det(\mathbf z_{\rm in},\mathbf z_{\rm out})/(\|\mathbf z_{\rm in}\|\|\mathbf z_{\rm out}\|)\);
\(\mathbf b\cdot\mathbf z_s=0\).

\begin{itemize}
\tightlist
\item
  \path|notes/chapters/ch05.md|
\item
  \path|notes/chapters/ch06.md|
\item
  \path|notes/chapters/ch10.md|
\item
  \path|tools/stellar_methods.py|
\item
  \path|calculation/src/adapters/ns_cell_shooting.f90|
\end{itemize}

\subsection{계산으로
옮기기}\label{uxacc4uxc0b0uxc73cuxb85c-uxc62euxae30uxae30}

제공 함수의 호출이 기본 경로다. 구간 보간·역함수·적응 적분·matching의
내부 코드를 작성하지 않고, 입력 조건과 반환된 근거를 해석한다.

\begin{enumerate}
\def\labelenumi{\arabic{enumi}.}
\tightlist
\item
  표면벡터와 determinant를 본문 식에서 설명한 뒤 제공 solve\_pair를
  호출한다.
\item
  run\_study의 단계별 결과에서 절대 주파수·차이·모드 대응을 각각 읽는다.
\item
  중심·표면·연결점·질량·무변형 및 적분 검사의 근거로 결론을 쓴다. 중심
  다음 항은 선택한다.
\end{enumerate}

\newpage

\subsection{실행해서 무엇을 볼
것인가}\label{uxc2e4uxd589uxd574uxc11c-uxbb34uxc5c7uxc744-uxbcfc-uxac83uxc778uxac00}

첫 명령 묶음은 한 pair의 시험이다. 전체 판정은 아래
\texttt{run\_study}의 모든 필요한 실행과 검사가 끝난 뒤에만 읽는다.

\begin{verbatim}
python tools/run_assignment.py A21 --plan
python tools/run_assignment.py A21 --run --output ../A21.json
\end{verbatim}

\begin{verbatim}
from tools.stellar_methods import run_study
r = run_study("../my_study", max_new_runs=2)
# 이후 같은 인수로 반복하거나 max_new_runs를 생략한다.
print(r["status"], len(r["pending_runs"]))
\end{verbatim}

\subsection{결과를
설명하기}\label{uxacb0uxacfcuxb97c-uxc124uxba85uxd558uxae30}

단일 pair 성공, 절대 주파수 수렴, 모드 대응과 전체 수치 여유는 별도
증거다. 적분 검사에서도 유한 표면항을 보존하고 저장 주파수를 다시 넣는
순환 검사를 피한다.

\textbf{남길 증거.} 같은 물리 BC와 두 방향 해의 연결; 두 별의 절대
주파수와 paired difference 수렴표; 고유함수 대응과 오차 항목의 계산
근거; study.json과 재개 가능한 실행 명령

\textbf{판단 기준.} 2점: 자동 pair 성공과 전체 numerical allowance를
구분하고, 같은 EOS·질량·mode 조건에서 결과의 부호를 방어한다. 최초
예측의 정오보다, 결과를 본 뒤 이유를 고치고 다른 입력에도 적용하는지를
본다.

\subsection{막혔을 때}\label{uxb9c9uxd614uxc744-uxb54c}

\textbf{힌트 1.} 두 비영 벡터가 평행하면 비자명한 선형해의 전체 배율을
맞출 수 있다.

\textbf{힌트 2.} 주파수만 비교하지 말고 matching점의 r,m와 고유함수도
확인한다.

\subsection{다음 만남에 다시
묻기}\label{uxb2e4uxc74c-uxb9ccuxb0a8uxc5d0-uxb2e4uxc2dc-uxbb3buxae30}

다음 만남에는 결과의 마지막 자릿수를 가린다. 같은 물리 문제인지 확인할
조건과 가장 취약한 오차원에 대한 다음 실험 하나를 스스로 고른다. 먼저
짧게 답한 뒤 본문과 자신의 기록을 열어 확인한다.

\textbf{범위.} Empirical allowance는 interval proof나 관측 신뢰구간이
아니다.

\end{document}
